\documentclass{article}
\usepackage[margin=1in]{geometry}
\usepackage{amsmath}
\usepackage{amssymb}
\usepackage{array}
\usepackage{graphicx}
\usepackage{booktabs}
\usepackage{hyperref}
\usepackage{xcolor}
\usepackage{enumitem}
\usepackage{setspace}
\usepackage{caption}
\usepackage{subcaption}
\usepackage{float}
\usepackage{parskip}
\usepackage{titlesec}
\usepackage{fancyhdr}
\usepackage{microtype}
\usepackage{multirow}

\pagestyle{fancy}
\fancyhf{}
\rhead{CS6370: Natural Language Processing}
\lhead{Information Retrieval Project}
\cfoot{\thepage}

\hypersetup{
    colorlinks=true,
    linkcolor=blue!70!black,
    urlcolor=blue!70!black,
    citecolor=blue!70!black
}

\titleformat{\section}{\large\bfseries}{\thesection.}{0.5em}{}[\titlerule]
\titleformat{\subsection}{\normalsize\bfseries}{\thesubsection.}{0.5em}{}
\titleformat{\subsubsection}{\normalsize\itshape}{\thesubsubsection.}{0.5em}{}

\title{
    \vspace{-1cm}
    \textbf{CS6370: Natural Language Processing}\\[0.3em]
    \large Project: Information Retrieval System Design, Evaluation, and Improvement\\[0.3em]
    \normalsize Cranfield Aerodynamics Dataset
}
\author{Student Report}
\date{}

\begin{document}

\maketitle
\thispagestyle{fancy}

\tableofcontents
\newpage

% ===========================================================================
\section{Part 1: Working out a Toy IR System}
% ===========================================================================

\subsection{Preprocessing and Inverted Index}

The three source documents are as follows: $d_1$: \textit{``The star in our solar system provides heat and light.''}, $d_2$: \textit{``That Hollywood star walked the red carpet for the movie premiere.''}, and $d_3$: \textit{``Astronomers observe distant stars and galaxies using telescopes.''}

\textbf{Step 1 — Tokenisation.} Each document is converted to lowercase and punctuation is stripped, producing:
\begin{itemize}[noitemsep]
    \item $d_1$: \texttt{[the, star, in, our, solar, system, provides, heat, and, light]}
    \item $d_2$: \texttt{[that, hollywood, star, walked, the, red, carpet, for, the, movie, premiere]}
    \item $d_3$: \texttt{[astronomers, observe, distant, stars, and, galaxies, using, telescopes]}
\end{itemize}

\textbf{Step 2 — Stopword removal.} Removing $\{\text{the, in, our, and, that, for}\}$ gives:
\begin{itemize}[noitemsep]
    \item $d_1$: \texttt{[star, solar, system, provides, heat, light]}
    \item $d_2$: \texttt{[hollywood, star, walked, red, carpet, movie, premiere]}
    \item $d_3$: \texttt{[astronomers, observe, distant, stars, galaxies, using, telescopes]}
\end{itemize}
Note that \textit{stars} in $d_3$ is a distinct type from \textit{star}, since no stemming was applied.

\textbf{Step 3 — Inverted Index.}

\begin{center}
\renewcommand{\arraystretch}{1.15}
\begin{tabular}{ll|ll}
\toprule
\textbf{Term} & \textbf{Postings} & \textbf{Term} & \textbf{Postings} \\
\midrule
astronomers & $\{d_3\}$ & observe   & $\{d_3\}$ \\
carpet      & $\{d_2\}$ & premiere  & $\{d_2\}$ \\
distant     & $\{d_3\}$ & provides  & $\{d_1\}$ \\
galaxies    & $\{d_3\}$ & red       & $\{d_2\}$ \\
heat        & $\{d_1\}$ & solar     & $\{d_1\}$ \\
hollywood   & $\{d_2\}$ & \textbf{star} & $\mathbf{\{d_1, d_2\}}$ \\
light       & $\{d_1\}$ & stars     & $\{d_3\}$ \\
movie       & $\{d_2\}$ & system    & $\{d_1\}$ \\
&                       & telescopes& $\{d_3\}$ \\
&                       & using     & $\{d_3\}$ \\
&                       & walked    & $\{d_2\}$ \\
\bottomrule
\end{tabular}
\end{center}

\subsection{TF-IDF Term-Document Matrix}

With $N=3$ documents, $\mathrm{IDF}(t) = \log_{10}(N / \mathrm{df}(t))$. Since all terms have $\mathrm{df}=1$ except \textit{star} with $\mathrm{df}=2$, we get $\mathrm{IDF}(\text{star}) = \log_{10}(1.5) \approx 0.176$ and $\mathrm{IDF}(\text{other}) = \log_{10}(3) \approx 0.477$. Each term has $\mathrm{TF}=1$ wherever it appears, making the final TF-IDF weight equal to IDF in that document.

\begin{table}[H]
\centering
\renewcommand{\arraystretch}{1.2}
\begin{tabular}{lccccc}
\toprule
\textbf{Term} & \textbf{df} & \textbf{IDF} & $w(d_1)$ & $w(d_2)$ & $w(d_3)$ \\
\midrule
heat       & 1 & 0.477 & 0.477 & 0     & 0     \\
hollywood  & 1 & 0.477 & 0     & 0.477 & 0     \\
light      & 1 & 0.477 & 0.477 & 0     & 0     \\
movie      & 1 & 0.477 & 0     & 0.477 & 0     \\
provides   & 1 & 0.477 & 0.477 & 0     & 0     \\
solar      & 1 & 0.477 & 0.477 & 0     & 0     \\
\textbf{star} & \textbf{2} & \textbf{0.176} & \textbf{0.176} & \textbf{0.176} & 0 \\
system     & 1 & 0.477 & 0.477 & 0     & 0     \\
stars/others & 1 & 0.477 & 0 & 0 & 0.477 \\
\bottomrule
\end{tabular}
\caption{Selected TF-IDF weights (19 unique terms total; 0 entries omitted for clarity).}
\end{table}

\subsection{Boolean Retrieval}

Query: \textit{``star light''}. Looking up the inverted index: \texttt{star} $\to \{d_1, d_2\}$, \texttt{light} $\to \{d_1\}$. Boolean AND (intersection) yields $\mathbf{d_1}$ as the sole matching document. With Boolean OR we would retrieve $d_1 \cup d_2$, but AND is the standard default for multi-term queries.

\subsection{Cosine Similarity and Ranking}

The query vector $\mathbf{q}$ has non-zero entries only for \textit{star} ($0.176$) and \textit{light} ($0.477$), giving $\|\mathbf{q}\| = \sqrt{0.176^2 + 0.477^2} \approx 0.508$.

\[
\mathrm{Cosine}(q, d_1) = \frac{(0.176)^2 + (0.477)^2}{0.508 \times \sqrt{0.176^2 + 5 \times 0.477^2}} = \frac{0.2585}{0.508 \times 1.081} \approx \mathbf{0.470}
\]
\[
\mathrm{Cosine}(q, d_2) = \frac{(0.176)^2}{0.508 \times \sqrt{0.176^2 + 6 \times 0.477^2}} \approx \frac{0.031}{0.600} \approx \mathbf{0.051}
\]
\[
\mathrm{Cosine}(q, d_3) = 0 \quad \text{(no shared terms)}
\]

\begin{itemize}[noitemsep]
    \item \textbf{Final ranking:} $d_1 \succ d_2 \succ d_3$, with cosine similarities 0.470, 0.051, and 0 respectively.
    \item \textbf{Is the ranking desirable?} Yes. The query \textit{``star light''} is about starlight as a physical phenomenon. $d_1$ discusses the Sun providing heat and light --- a direct match. $d_2$ only shares the token \textit{star} used in the celebrity sense, so its score (0.051) is much lower, and it is correctly demoted. $d_3$ scores zero because it has no query terms at all.
\end{itemize}

\subsection{Word Sense Ambiguity}

Query: \textit{movie star''}.

\begin{itemize}[noitemsep]
    \item \textbf{Which document(s) should ideally be retrieved?} Only $d_2$, which describes a Hollywood film premiere --- the celebrity'' sense of \textit{star} matches exactly. $d_1$ and $d_3$ are irrelevant.

    \item \textbf{Does the system retrieve the correct document(s)?} Not exclusively. After stopword removal the query terms are \textit{movie} (IDF = 0.477) and \textit{star} (IDF = 0.176). Computing cosine similarity gives $d_2 \approx 0.431$ and $d_1 \approx 0.056$ --- so $d_1$ is also returned with a non-zero score, even though it has nothing to do with celebrities. The system fails to restrict retrieval to the correct document.

    \item \textbf{How does word sense ambiguity affect retrieval?} The TF-IDF model assigns the same token \textit{star} to both $d_1$ (astronomical star) and $d_2$ (celebrity). Because both documents contain this token, both receive a non-zero score from the query. This is the classical \textbf{polysemy} problem: the bag-of-words representation cannot separate the two word senses, so the system returns an irrelevant document alongside the correct one, degrading precision.
\end{itemize}

% ===========================================================================
\section{Part 2: Building an IR System — Implementation}
% ===========================================================================

\subsection{TF-IDF Vector Space Model}

The retrieval pipeline was implemented in \texttt{informationRetrieval.py} following the standard Vector Space Model (VSM). The implementation uses an inverted index (a dictionary mapping every vocabulary term to its per-document raw term frequency) combined with augmented TF weighting and standard IDF:

\[
w(t, d) = \underbrace{\left(0.5 + 0.5 \cdot \frac{\mathrm{tf}(t,d)}{\max_{t'}\mathrm{tf}(t',d)}\right)}_{\text{augmented TF}} \times \underbrace{\log_{10}\!\left(\frac{N}{\mathrm{df}(t)}\right)}_{\text{IDF}}
\]

Augmented TF prevents longer documents from dominating by normalising raw term counts against the most frequent term in the same document. Documents are ranked by cosine similarity in the resulting vector space. The implementation handles out-of-vocabulary (OOV) terms gracefully: when a query term is unseen, its IDF lookup returns zero and the term contributes nothing to the dot product.

The same preprocessing pipeline was applied uniformly to both queries and documents before indexing:
\begin{itemize}[noitemsep]
    \item Punkt sentence segmentation
    \item Penn Treebank tokenisation
    \item WordNet lemmatisation
    \item NLTK stopword removal
\end{itemize}

% ===========================================================================
\section{Part 3: Evaluating the IR System}
% ===========================================================================

\subsection{Evaluation Measures}

Six evaluation measures were implemented in \texttt{evaluation.py}:

\textbf{Precision@k} counts the fraction of retrieved documents in the top-$k$ that are truly relevant:
$\mathrm{Precision@k} = |\{d \in R_k \mid d \in \mathcal{R}\}| / k$.

\textbf{Recall@k} measures what proportion of all relevant documents were retrieved in the top-$k$:
$\mathrm{Recall@k} = |\{d \in R_k \mid d \in \mathcal{R}\}| / |\mathcal{R}|$.

\textbf{F$_{0.5}$-score@k} uses the general $F_\beta$ harmonic mean formula:
$F_\beta = \frac{(1+\beta^2) \cdot P \cdot R}{\beta^2 \cdot P + R}$.
Substituting $\beta = 0.5$ ($\beta^2 = 0.25$) gives $F_{0.5} = \frac{1.25 \cdot P \cdot R}{0.25P + R}$, which weights precision twice as heavily as recall.

\textbf{Average Precision (AP@k)} rewards systems that place relevant documents early in the ranking. For a query with $|\mathcal{R}|$ total relevant documents,
$\mathrm{AP@k} = \frac{1}{|\mathcal{R}|} \sum_{i=1}^{k} \mathrm{Precision@i} \cdot \mathbf{1}[\mathrm{doc}_i \in \mathcal{R}]$,
where $\mathbf{1}[\mathrm{doc}_i \in \mathcal{R}] = 1$ if the document at rank $i$ is relevant and 0 otherwise --- so only ranks that actually hit a relevant document contribute to the sum.

\textbf{nDCG@k} normalises the Discounted Cumulative Gain against the ideal ordering. Unlike precision-based metrics, DCG is designed for \emph{graded} relevance: the numerator $\mathrm{rel}_i$ is the relevance grade of the document retrieved at rank $i$ (for Cranfield, this is the qrel position value 1--4, or 0 if not in the qrels), and the denominator $\log_2(i+1)$ is the positional discount --- a relevant document at rank 1 contributes more than the same document at rank 5. $\mathrm{DCG@k} = \sum_{i=1}^{k} \frac{\mathrm{rel}_i}{\log_2(i+1)}$, and $\mathrm{nDCG@k} = \mathrm{DCG@k} / \mathrm{IDCG@k}$. In our main evaluation we use the binary simplification $\mathrm{rel}_i \in \{0,1\}$ (qrel positions 1--4 all treated as equally relevant, as required by the assignment); the graded variant nDCG@10$^*$ restores grades 2 and 1 for qrel positions 1--2 and 3--4 respectively.

\textbf{Mean Reciprocal Rank (MRR@k)} measures how early the first relevant document appears. For each query the reciprocal rank is $1/r$ where $r$ is the position (within the top-$k$) of the first relevant document, or 0 if none is found. $\mathrm{MRR@k} = \frac{1}{|Q|} \sum_{q} \frac{1}{r_q}$.

Relevance was determined by the Cranfield qrels: documents with a position value of 1, 2, 3, or 4 were treated as relevant.

\subsection{Quantitative Results}

\begin{table}[H]
\centering
\renewcommand{\arraystretch}{1.25}
\resizebox{\linewidth}{!}{%
\begin{tabular}{lcccccc}
\toprule
\textbf{k} & \textbf{Precision} & \textbf{Recall} & \textbf{F$_{0.5}$} & \textbf{MAP} & \textbf{nDCG} & \textbf{MRR} \\
\midrule
1  & 0.6000 & 0.1048 & 0.2872 & 0.1048 & 0.6000 & 0.6000 \\
2  & 0.4800 & 0.1616 & 0.3195 & 0.1540 & 0.5072 & 0.6533 \\
3  & 0.4311 & 0.2115 & 0.3328 & 0.1894 & 0.4690 & 0.6800 \\
4  & 0.3922 & 0.2464 & 0.3287 & 0.2115 & 0.4482 & 0.6889 \\
5  & 0.3591 & 0.2789 & 0.3199 & 0.2297 & 0.4343 & 0.6933 \\
6  & 0.3311 & 0.3044 & 0.3075 & 0.2427 & 0.4266 & 0.6963 \\
7  & 0.3079 & 0.3247 & 0.2951 & 0.2522 & 0.4222 & 0.6976 \\
8  & 0.2867 & 0.3445 & 0.2822 & 0.2588 & 0.4210 & 0.7009 \\
9  & 0.2672 & 0.3589 & 0.2688 & 0.2643 & 0.4201 & 0.7019 \\
\textbf{10} & \textbf{0.2507} & \textbf{0.3710} & \textbf{0.2568} & \textbf{0.2691} & \textbf{0.4203} & \textbf{0.7028} \\
\bottomrule
\end{tabular}}
\caption{VSM Baseline evaluation metrics averaged over all 225 Cranfield queries.}
\label{tab:baseline}
\end{table}

\subsection{Global Metrics}

\textbf{MAP@10: 0.2691.} \quad \textbf{MRR@10: 0.7028.}
Mean Reciprocal Rank at cutoff $k$ gives the average, over all 225 queries, of $1/r$ where $r$ is the rank of the first relevant document found within the top-$k$ results (0 if none appears). At $k{=}1$, MRR collapses to Precision@1 = 0.6000 because the only available rank is 1. As $k$ grows, queries whose first relevant document sits at position 2 or 3 start contributing a non-zero reciprocal rank, so MRR rises steadily: 0.6533 at $k{=}2$, 0.6800 at $k{=}3$, reaching 0.7028 at $k{=}10$. MRR@10 = 0.7028 tells us that on most queries the first relevant document was found at rank 1 or 2. For the queries where rank 1 was wrong, the system still managed to get a relevant document into rank 2 or 3 in most cases, which is why MRR is noticeably higher than Precision@1.

\subsection{Evaluation Plots}

\begin{figure}[H]
\centering
\includegraphics[width=0.82\textwidth]{eval_plot_part3.png}
\caption{VSM baseline: Precision, Recall, F$_{0.5}$, MAP, nDCG, and MRR plotted against cutoff $k$ over all 225 Cranfield queries.}
\label{fig:part3}
\end{figure}

\subsection{Observations from the Plots}

\textbf{Precision vs.\ $k$} declines monotonically from 0.60 at $k{=}1$ to 0.25 at $k{=}10$. This is the standard precision--recall trade-off: expanding the result window retrieves more relevant documents but necessarily also more irrelevant ones, diluting precision.

\textbf{Recall vs.\ $k$} increases monotonically, from 0.105 to 0.371. Because the Cranfield queries have an average of roughly 9 relevant documents each, a cutoff of $k{=}10$ captures just over a third of all relevant material --- indicating substantial room for improvement.

\textbf{F$_{0.5}$-score} peaks at $k{=}3$ (0.333) and then declines, reflecting the higher weight placed on precision. Once precision degrades past a certain point, the combined score falls even though recall continues to improve.

\textbf{MAP vs.\ $k$} rises gradually from 0.10 at $k{=}1$ to 0.27 at $k{=}10$. Each time a relevant document is retrieved at a new position, it adds a precision-at-that-position value to the AP sum, so MAP can only go up or stay flat as $k$ grows.

\textbf{nDCG vs.\ $k$} moves in the \emph{opposite} direction: it decreases from 0.60 at $k{=}1$ to approximately 0.42 at $k{=}10$, then plateaus. At $k{=}1$ there is only one rank position, so nDCG collapses to Precision@1 = 0.60 (the system either placed a relevant document first or it did not). As $k$ grows, additional non-relevant documents accumulate in the ranked list, adding nothing to the DCG numerator while the ideal DCG denominator continues to grow, so the ratio falls. The plateau near 0.42 past $k{=}5$ shows diminishing change --- most of the useful signal is already captured in the top five positions.

\subsection{Run Time}

The complete pipeline (loading, preprocessing 1400 documents and 225 queries, building the inverted index, ranking all queries, and computing all six evaluation metrics) completed in approximately \textbf{4.4 seconds}, demonstrating the efficiency of the sparse inverted-index approach.

% ===========================================================================
\section{Part 4: Analysis --- Limitations of the VSM}
% ===========================================================================

\subsection{Vocabulary Mismatch and the Assumption of Orthogonality}

The fundamental assumption of the Vector Space Model is that every distinct word type is a completely independent, orthogonal dimension of meaning. This is wrong in general, and catastrophically wrong for scientific text. Consider the query \textit{``what similarity laws must be obeyed when constructing aeroelastic models of heated high speed aircraft''} and Document 29, which discusses \textit{``transient temperatures and thermal stresses in simple models intended to simulate parts of an aircraft structure subjected to aerodynamic heating.''}

The qrels mark Document 29 as highly relevant (position 2). Yet our TF-IDF system placed it at \textbf{rank 109}. The cause is a direct vocabulary mismatch: the query uses \textit{heated}, while Document 29 uses \textit{thermal}, \textit{temperatures}, and \textit{heating}. Since these are distinct word types in the vector space, the dot product receives no contribution from this conceptual overlap, and the cosine similarity collapses to near zero.

This is not an edge case. The Cranfield corpus is composed of highly technical aerodynamics abstracts, where synonymy between specialist terms (e.g., \textit{velocity} vs.\ \textit{speed}, \textit{pressure} vs.\ \textit{stress}, \textit{viscous} vs.\ \textit{viscid}) is pervasive. The Bag-of-Words assumption is particularly damaging in this domain.

\subsection{The Out-of-Vocabulary (OOV) Problem}

When a query contains no terms that appear in the document corpus, the query vector reduces to the zero vector. Every document then scores a cosine similarity of exactly zero, and the ranking degenerates to whichever arbitrary order the documents were loaded in. In practice, the system would return $[1, 2, 3, 4, 5, \ldots]$ --- a meaningless artefact of file order.

A controlled test confirmed this behaviour: querying \textit{``intergalactic starships flippity floppity''} returned the top-5 documents as \texttt{[1, 2, 3, 4, 5]}, confirming the total absence of discriminative information in the response. For a mixed query like \textit{``aerodynamics of intergalactic starships''}, only the known term \textit{aerodynamics} contributed to the ranking, effectively penalising the user for any novel terminology in their information need.

\subsection{Additional Limitations}

Beyond vocabulary mismatch and OOV, the VSM has three further structural problems:

\begin{itemize}[noitemsep]
    \item \textbf{Word order blindness.} The model treats \textit{boundary layer} and \textit{layer boundary} as identical because it only counts tokens, not their positions. Document structure (titles, headings, abstracts) is also completely discarded.
    \item \textbf{Length bias.} Even with augmented TF normalisation, a document twice as long still has roughly twice the opportunity to accumulate term matches. So a long document with moderate relevance can outscore a short, tightly-focused one just because it has more words.
    \item \textbf{Negation blindness.} The query \textit{``without turbulence''} produces the same token set as \textit{``with turbulence''}, so both produce the same query vector. There is simply no way in this model to say one word cancels out another.
\end{itemize}

% ===========================================================================
\section{Part 5: Improving the IR System}
% ===========================================================================

\subsection{Introduction and Problem Framing}

The failure analysis in Part 4 shows a clear and measurable problem: the VSM gets MAP@10 of only 0.2691 on Cranfield, and for Query 1, Document 29 was placed at rank 109 even though the qrels mark it as highly relevant. The root cause is the orthogonality assumption, where every word type gets its own independent dimension regardless of meaning overlap. Our working hypothesis is that if we can somehow make the retrieval model aware of synonymy or co-occurrence relationships between terms, it should be able to fix at least some of these failures.

To check this, we implemented and evaluated six alternative retrieval systems against the same 225 queries and 1400 documents. For deciding whether any improvement is actually real and not just noise, we use the Wilcoxon Signed-Rank Test on per-query Average Precision values. AP is chosen as the per-query unit because it summarises the full ranked list for one query in a single number --- it rewards finding relevant documents and finding them early, making it more informative than a single cutoff metric like Precision@10. It also produces one value per query per system, which is exactly what a paired test needs: 225 matched pairs, one per query. For methods where theory predicts a clear improvement (BM25, LSA, PRF), we use a one-sided test. For methods whose direction of effect is uncertain (query expansion, ESA), we use a two-sided test, which is the statistically correct choice when no directional prior exists.

\subsection{Proposed Approaches}

\subsubsection{Method A --- BM25 (Okapi BM25)}

\textit{\textbf{Formulation:} Took Okapi BM25 over Cranfield (1400 documents, 225 queries) under the assumption that TF saturation and length normalisation better model term informativeness than augmented TF-IDF, across Precision@k, Recall@k, MAP@k, nDCG@k (binary and graded), and MRR@k. Hyperparameters $k_1 \in \{1.2, 1.5, 2.0\}$ and $b \in \{0.5, 0.75, 0.9\}$ selected by 5-fold cross-validation.}

BM25 replaces augmented TF with a saturating TF weighting and a length-normalised document model, addressing two systematic weaknesses of TF-IDF. The scoring function for a term $t$ in document $d$ is:

\[
\mathrm{score}(d, Q) = \sum_{t \in Q} \mathrm{IDF}^*(t) \cdot \frac{\mathrm{tf}(t,d) \cdot (k_1 + 1)}{\mathrm{tf}(t,d) + k_1\!\left(1 - b + b \cdot \frac{|d|}{\mathrm{avgdl}}\right)}
\]

where $\mathrm{IDF}^*(t) = \log\!\left(\frac{N - \mathrm{df}(t) + 0.5}{\mathrm{df}(t) + 0.5} + 1\right)$ (Robertson IDF, always positive). The parameter $k_1$ controls TF saturation (a term appearing 100 times in a document is not 100$\times$ more relevant than one appearing once) and $b$ controls the degree of length normalisation. Hyperparameters $k_1$ and $b$ were selected by 5-fold cross-validation on query-level MAP, with the grid $k_1 \in \{1.2, 1.5, 2.0\}$ and $b \in \{0.5, 0.75, 0.9\}$. The optimal configuration found was $k_1{=}2.0$, $b{=}0.9$, indicating that on this aerodynamics corpus heavier TF saturation and stronger length normalisation are beneficial, likely because the document collection has high variance in abstract length.

\subsubsection{Method B --- Latent Semantic Analysis (LSA)}

\textit{\textbf{Formulation:} Took LSA (Truncated SVD) over Cranfield under the assumption that the $k$-rank approximation of the TF-IDF term-document matrix captures latent semantic structure that reduces vocabulary mismatch, across the same measures. The number of latent dimensions selected from $\{50, 100, 200, 300\}$ by 5-fold cross-validation.}

LSA decomposes the TF-IDF term-document matrix $M \in \mathbb{R}^{V \times N}$ using Truncated Singular Value Decomposition into $M \approx U_k \Sigma_k V_k^\top$. Each document is represented by its corresponding \emph{column} in $\Sigma_k V_k^\top$ (a $k$-dimensional latent vector). A query is projected into the same space by the standard folding-in formula $\mathbf{q}_{\mathrm{lsa}} = \Sigma_k^{-1} U_k^\top \mathbf{q}_{\mathrm{tfidf}}$: $U_k^\top$ maps from the $V$-dimensional vocabulary into concept space, and $\Sigma_k^{-1}$ places the query on the same scale as the document vectors. In practice $\Sigma_k^{-1}$ is dropped for cosine similarity since it rescales the query uniformly and does not change rankings.

In the $k$-dimensional latent space, terms that consistently co-occur in similar documents are pushed close together even if they never literally co-occur in the same sentence. This directly addresses the vocabulary mismatch problem: \textit{heated} and \textit{thermal} will have nearby latent representations because they appear in documents that share many other technical terms.

The number of latent components $k$ is a critical hyperparameter. We ran 5-fold cross-validation over $k \in \{50, 100, 200, 300\}$ on query-level MAP to select the best $k$ without leaking test-set information. The CV selected $k{=}200$. We also ran a full ablation on all 225 queries for the purpose of plotting the MAP-vs-$k$ curve, shown below:

\begin{table}[H]
\centering
\renewcommand{\arraystretch}{1.2}
\begin{tabular}{lcccc}
\toprule
\textbf{Components ($k$)} & \textbf{Precision@10} & \textbf{Recall@10} & \textbf{MAP@10} & \textbf{nDCG@10} \\
\midrule
50  & 0.2667 & 0.3787 & 0.2553 & 0.4024 \\
100 & 0.2813 & 0.4021 & 0.2882 & 0.4386 \\
\textbf{200} & \textbf{0.3004} & \textbf{0.4347} & \textbf{0.3233} & \textbf{0.4804} \\
300 & 0.2973 & 0.4304 & 0.3205 & 0.4795 \\
VSM Baseline & 0.2507 & 0.3710 & 0.2691 & 0.4203 \\
\bottomrule
\end{tabular}
\caption{LSA ablation over all 225 queries. MAP@10 peaks at $k{=}200$ and declines at $k{=}300$, consistent with the bias-variance trade-off. 5-fold CV independently confirmed $k{=}200$ as the best choice (CV MAP@10 estimate = 0.3233). Note: the ablation plot is for illustration; the CV-selected $k$ is what was used in the final comparison.}
\label{tab:ablation}
\end{table}

\begin{figure}[H]
\centering
\includegraphics[width=0.82\textwidth]{part5_lsa_ablation.png}
\caption{LSA ablation: MAP@10 (left) and nDCG@10 (right) as a function of the number of latent dimensions $k$. The red dashed line is the VSM baseline. Performance peaks at $k{=}200$ and the 5-fold CV independently selected this value.}
\label{fig:lsa_ablation}
\end{figure}

\subsubsection{Method C --- VSM with WordNet Query Expansion}

\textit{\textbf{Formulation:} Took VSM with WordNet Query Expansion over Cranfield under the assumption that first-synset WordNet synonyms (at most 3 per query term) bridge vocabulary gaps without introducing excessive query drift, across the same measures.}

Query Expansion (QE) attempts to address vocabulary mismatch at query time rather than corpus time. For each token in the preprocessed query, synonyms are retrieved from the primary WordNet synset (at most 3 per word, to avoid query drift). These synonyms are appended as additional query tokens and scored through the existing TF-IDF index.

The key risk is \emph{query drift}: if the first synset is ambiguous (e.g., \textit{star} retrieving astronomy synonyms for a query about celebrities), the expanded terms pollute the result. This risk is especially acute for short, ambiguous queries.

\subsubsection{Method D --- BM25 with Pseudo-Relevance Feedback (PRF)}

\textit{\textbf{Formulation:} Took BM25 with Pseudo-Relevance Feedback (top-3 first-pass documents assumed relevant, top-3 new terms injected into a second BM25 pass) over Cranfield under the assumption that the first-pass top documents are sufficiently on-topic to provide useful expansion vocabulary, across the same measures.}

PRF combines BM25 for initial retrieval with an automatic query expansion step. The top 3 documents from the first BM25 pass are assumed relevant, and the 3 most frequent new terms from those documents are injected into a second BM25 pass. This allows the system to adapt vocabulary to the domain of the query --- for instance, a query about \textit{boundary layers} might expand to include co-occurring terms like \textit{laminar}, \textit{turbulent}, and \textit{separation}.

The central limitation is the same as for QE: if the initial top-3 documents are irrelevant, the feedback terms will actively mislead the second pass (query drift). Our failure analysis confirmed this for queries in low-density relevance regions of the corpus.

\subsubsection{Method E and F --- Explicit Semantic Analysis (ESA)}

\textit{\textbf{Formulation (Method E):} Took ESA with WordNet concept space over Cranfield under the assumption that WordNet synsets form a useful concept basis for aerodynamics text, with document/query vectors built by TF-IDF-weighted synset activations, across the same measures.}

\textit{\textbf{Formulation (Method F):} Took ESA with Cranfield concept space over Cranfield under the assumption that the 1400 Cranfield documents themselves form a coherent concept basis (each document treated as a concept dimension), which avoids domain mismatch, across the same measures.}

ESA projects both documents and queries into a high-dimensional \emph{concept space} rather than word space. We tested two concept spaces: (E) a WordNet synset space, where each synset is a concept dimension, and (F) a Cranfield-internal concept space, where each document is itself treated as a concept (the document-level ESA variant). In both cases, a document's weight for concept $c$ accumulates the TF-IDF weight of every term in the document that maps to $c$. Ranking is by cosine similarity in the concept space.

The intuition is that two documents describing similar phenomena will activate similar sets of concepts even if they use different surface vocabulary.

\subsection{Experimental Setup}

All methods were evaluated on the full Cranfield test collection: 1400 documents and 225 queries with qrels defining binary relevance at positions 1--4. Metrics are computed at cutoffs $k{=}1$ to $k{=}10$. No information from qrels was used during system construction (no supervised feedback). Both BM25 hyperparameters ($k_1$, $b$) and LSA's $n\_components$ were each selected by independent 5-fold cross-validation on the query set, with no access to qrel labels during CV. We also report MRR@10 (Mean Reciprocal Rank), which is particularly informative for navigational queries where the first relevant document matters most, and graded nDCG@10 where position~1--2 qrel entries contribute gain~2 and position~3--4 entries contribute gain~1.

\subsection{Quantitative Results}

\begin{table}[H]
\centering
\renewcommand{\arraystretch}{1.3}
\begin{tabular}{lcccccc}
\toprule
\textbf{System} & \textbf{Prec@10} & \textbf{Rec@10} & \textbf{MAP@10} & \textbf{nDCG@10} & \textbf{nDCG@10$^*$} & \textbf{MRR@10} \\
\midrule
VSM (Baseline) & 0.2507 & 0.3710 & 0.2691 & 0.4203 & 0.4249 & 0.7028 \\
\midrule
BM25 ($k_1{=}2.0$, $b{=}0.9$) & 0.2898 & 0.4218 & \textbf{0.3304} & 0.4918 & 0.4933 & \textbf{0.8002} \\
BM25 + PRF     & 0.3004 & 0.4261 & \textbf{0.3388} & \textbf{0.4914} & \textbf{0.4887} & 0.7590 \\
LSA ($k{=}200$, CV)   & 0.3004 & 0.4347 & 0.3233 & 0.4804 & 0.4733 & 0.7249 \\
ESA (Cranfield)& 0.2773 & 0.3974 & 0.2963 & 0.4516 & 0.4506 & 0.7246 \\
VSM + QE       & 0.2156 & 0.3301 & 0.2299 & 0.3647 & 0.3726 & 0.6249 \\
ESA (WordNet)  & 0.2076 & 0.3059 & 0.2115 & 0.3478 & 0.3563 & 0.6157 \\
\bottomrule
\end{tabular}
\caption{Comparative evaluation at $k{=}10$. nDCG@10$^*$ is graded nDCG using gains 2 (positions 1--2 in qrels) and 1 (positions 3--4). Rows below the VSM are sorted by MAP@10. Bold entries are the two best values per column.}
\label{tab:comparison}
\end{table}

\begin{figure}[H]
\centering
\begin{subfigure}[b]{0.99\textwidth}
    \centering
    \includegraphics[width=\textwidth]{part5_comparison_map_ndcg.png}
    \caption{MAP@k (left) and nDCG@k (right) vs.\ cutoff $k$ for all seven systems.}
    \label{fig:map_ndcg}
\end{subfigure}
\vspace{0.5em}
\begin{subfigure}[b]{0.99\textwidth}
    \centering
    \includegraphics[width=\textwidth]{part5_final_bar_chart.png}
    \caption{Bar chart of Precision, Recall, MAP, nDCG, and MRR at $k{=}10$ for all systems.}
    \label{fig:bar}
\end{subfigure}
\caption{Comparative retrieval performance across all seven systems on the Cranfield dataset.}
\label{fig:comparison}
\end{figure}

\subsection{Statistical Significance Testing}

We use the Wilcoxon Signed-Rank Test (non-parametric; no Gaussian assumption required) on per-query AP@10 distributions, with significance level $\alpha = 0.05$. Crucially, we apply a \textbf{one-sided} test ($H_1$: method $>$ baseline) only for methods where theory provides a strong directional prior (BM25, LSA, PRF). For methods where the direction of change is not guaranteed (VSM+QE, ESA), we apply a \textbf{two-sided} test, which is statistically correct when no directional assumption is justified. Using one-sided tests for all methods, as is sometimes done, would mask the significance of harm from degrading methods.

\textbf{Formal hypothesis framing:}

\begin{itemize}[noitemsep]
    \item $H_0$: The per-query AP scores of method $X$ and the VSM baseline come from the same distribution.
    \item $H_1^{>}$ (one-sided): Method $X$ produces \emph{strictly higher} AP scores (BM25, LSA, PRF).
    \item $H_1^{\neq}$ (two-sided): Method $X$ produces \emph{significantly different} AP scores (QE, ESA).
\end{itemize}

Effect sizes are measured by Cohen's $d = \bar{\delta} / \sigma_\delta$ on the paired differences, where $\bar{\delta}$ is the mean difference and $\sigma_\delta$ is the standard deviation of differences.

\begin{table}[H]
\centering
\renewcommand{\arraystretch}{1.3}
\begin{tabular}{llccccl}
\toprule
\textbf{System} & \textbf{Test} & \textbf{$W$} & \textbf{$p$-value} & \textbf{$\bar{\delta}$ MAP} & \textbf{Cohen's $d$} & \textbf{Decision} \\
\midrule
BM25           & one-sided & 13663.5 & $3.3 \times 10^{-16}$ & $+0.061$ & $+0.569$ (med) & \textbf{Reject $H_0$} \\
BM25 + PRF     & one-sided & 14079.0 & $1.8 \times 10^{-12}$ & $+0.070$ & $+0.494$ (small) & \textbf{Reject $H_0$} \\
LSA ($k{=}200$, CV) & one-sided & 13018.5 & $2.1 \times 10^{-9}$  & $+0.054$ & $+0.434$ (small) & \textbf{Reject $H_0$} \\
ESA (Cranfield)& two-sided & 6135.0  & $7.3 \times 10^{-5}$  & $+0.027$ & $+0.255$ (small) & \textbf{Reject $H_0$} \\
VSM + QE       & two-sided & 2237.5  & $8.1 \times 10^{-13}$ & $-0.039$ & $-0.415$ (small) & \textbf{Reject $H_0$ (\emph{sig.\ worse})} \\
ESA (WordNet)  & two-sided & 4082.0  & $4.0 \times 10^{-9}$  & $-0.058$ & $-0.383$ (small) & \textbf{Reject $H_0$ (\emph{sig.\ worse})} \\
\bottomrule
\end{tabular}
\caption{Wilcoxon Signed-Rank test results ($\alpha = 0.05$, $n{=}225$ queries). $\bar{\delta}$ MAP is the mean per-query AP difference relative to VSM baseline. Negative $\bar{\delta}$ and $d$ indicate the method is significantly \emph{worse} than baseline.}
\label{tab:wilcoxon}
\end{table}

All six methods are statistically significantly different from the VSM baseline. The key insight from the two-sided tests is that VSM+QE and ESA-WordNet are not merely indistinguishable from the baseline as previously reported with one-sided tests --- they are \emph{significantly worse} ($p < 10^{-8}$, $\bar{\delta} < 0$). The one-sided test had masked this because it was testing in the wrong direction. We therefore formally conclude:

\begin{quote}
\textit{``BM25 ($k_1{=}2.0$, $b{=}0.9$, CV-selected) is better than the TF-IDF VSM baseline with respect to MAP@10 on Cranfield under the assumption that saturating TF and length normalisation better model term informativeness than augmented TF-IDF ($p = 3.3 \times 10^{-16}$, Cohen's $d = 0.57$, medium effect).''}
\end{quote}

\subsection{MRR and Graded nDCG}

MRR@10 (Table~\ref{tab:comparison}, rightmost column) tells a complementary story: BM25 achieves MRR = 0.8002, meaning the first relevant document appears on average within the first 1.25 positions --- strong navigational quality. VSM+QE drops MRR to 0.6249 and ESA-WordNet to 0.6157, confirming that both expansion methods harm early precision. Graded nDCG (column nDCG@10$^*$) rewards systems that rank grade-2 relevant documents (positions 1--2 in qrels) above grade-1. BM25+PRF achieves the highest graded nDCG (0.4887), suggesting it is best at surfacing the most relevant documents at the very top of the ranked list.

\subsection{Marginal Difference Analysis}

Three system pairs exhibited MAP differences below 2\%, requiring deeper analysis before making superiority claims:

\begin{table}[H]
\centering
\renewcommand{\arraystretch}{1.2}
\begin{tabular}{llcccc}
\toprule
\textbf{System A} & \textbf{System B} & \textbf{$\Delta$MAP} & \textbf{A wins} & \textbf{B wins} & \textbf{Ties} \\
\midrule
BM25+PRF & BM25            & 0.0083 & 92 queries & 81 queries  & 52 \\
BM25     & LSA ($k{=}200$, CV) & 0.0071 & 105 queries & 83 queries & 37 \\
BM25+PRF & LSA ($k{=}200$, CV) & 0.0154 & 110 queries & 82 queries & 33 \\
VSM+QE   & ESA (WordNet)   & 0.0184 & 91 queries  & 89 queries  & 45 \\
\bottomrule
\end{tabular}
\caption{Query-level win/loss breakdown for system pairs with MAP difference $<2\%$. Note the last row: both QE and ESA-WordNet are below baseline, so ``wins'' here means beats the other degraded method, not beats the baseline.}
\end{table}

The 0.71\% MAP gap between BM25 and LSA translates to 105 vs.\ 83 per-query wins --- a clear directional signal, but the gap is small enough that it does not survive as a decisive superiority claim without pair-level hypothesis testing. BM25+PRF wins on 110 queries against LSA (82), giving it the strongest per-query record among the top three methods. All three (BM25, BM25+PRF, LSA) represent genuinely comparable retrieval quality in the practical range.

Looking at the PRF-vs-BM25 pair (0.83\% MAP gap) through the qrel density lens reveals \emph{why} this marginal difference exists. PRF wins on 43 of 67 dense-qrel queries (64.2\%) but only 5 of 35 sparse-qrel queries (14.3\%). This follows directly from what PRF assumes: the first-pass BM25 top-3 are relevant enough to yield useful expansion terms. For dense queries with $\ge 10$ relevant documents in the collection, this assumption usually holds --- BM25 reliably surfaces at least one relevant document, and the expansion terms genuinely reflect the query topic. For sparse queries (1--3 relevant documents), first-pass reliability is lower, expansion terms come from irrelevant documents, and the second pass is actively misled --- which is why plain BM25 wins 45.7\% of sparse-query matchups vs.\ PRF's 14.3\%. The 0.83\% aggregate MAP advantage for PRF is therefore a density-conditional effect: PRF is concretely better for dense queries, but slightly worse for sparse ones, and the Cranfield distribution (67 dense, 35 sparse) tips the average in PRF's favour.

The BM25-vs-LSA gap (0.71\% MAP) shows a different pattern. BM25 wins 53.7\% of dense-query matchups vs.\ LSA's 40.3\%, and 40.0\% vs.\ 25.7\% for sparse. Unlike the PRF-vs-BM25 relationship, BM25's advantage over LSA is present across \emph{all} density classes and is actually larger for dense queries. This means LSA's latent-space dimensionality reduction, while helpful for vocabulary matching, does not fully compensate for BM25's superior term weighting (Robertson IDF, TF saturation) even when many relevant documents are available.

\subsection{Per-Query Granularity Analysis}

Aggregate MAP numbers conceal substantial variability across queries. Table~\ref{tab:pq} shows the number of queries on which each method improved, degraded, or produced no change relative to the baseline.

\begin{table}[H]
\centering
\renewcommand{\arraystretch}{1.2}
\begin{tabular}{lccc}
\toprule
\textbf{System} & \textbf{Queries Improved} & \textbf{Queries Degraded} & \textbf{Unchanged} \\
\midrule
BM25           & 142 & 37  & 46 \\
BM25 + PRF     & 136 & 52  & 37 \\
LSA ($k{=}200$, CV) & 124 & 62  & 39 \\
ESA (Cranfield)& 120 & 71  & 34 \\
ESA (WordNet)  & 60  & 120 & 45 \\
VSM + QE       & 28  & 130 & 67 \\
\bottomrule
\end{tabular}
\caption{Per-query win/loss counts against the VSM baseline ($k{=}10$). Rows sorted by queries improved.}
\label{tab:pq}
\end{table}

BM25's per-query profile is the most favourable: it improves 142 queries and degrades only 37, a ratio of nearly 4:1 in favour of improvement. BM25+PRF is second at 136/52. In contrast, VSM+QE degrades 130 queries while improving only 28 --- a 1:4.6 ratio against improvement. This is a strong signal that unconstrained WordNet expansion introduces more noise than signal on this highly technical corpus, a finding now confirmed at the level of statistical significance by the two-sided Wilcoxon test.

\subsection{Query-Class-Level Analysis}

Aggregate MAP numbers treat all 225 queries as a flat pool. Here we break them down by two axes to see \emph{for which class of query} each method wins or loses.

\textbf{Classification scheme.} Queries are classified by (1) \textbf{length}: short ($\le 5$ words), medium (6--10 words), long ($\ge 11$ words); and (2) \textbf{qrel density}: sparse (1--3 relevant documents), medium (4--9), dense ($\ge 10$). Among 225 Cranfield queries, none has $\le 5$ words --- all fall into medium (32 queries) or long (193 queries). By density: 35 sparse, 123 medium, 67 dense.

\begin{table}[H]
\centering
\renewcommand{\arraystretch}{1.2}
\begin{tabular}{lcc}
\toprule
\textbf{System} & \textbf{Medium (6--10 words)} & \textbf{Long ($\ge$11 words)} \\
\midrule
VSM (Baseline)      & 0.2110 & 0.2787 \\
BM25                & 0.2751 & 0.3396 \\
BM25 + PRF          & \textbf{0.2831} & \textbf{0.3480} \\
LSA ($k{=}200$, CV) & 0.2735 & 0.3316 \\
ESA (Cranfield)     & 0.2505 & 0.3040 \\
VSM + QE            & 0.1964 & 0.2355 \\
ESA (WordNet)       & 0.1506 & 0.2216 \\
\bottomrule
\end{tabular}
\caption{Mean AP@10 by query length class. Sorted by long-query MAP. Bold entries are best per column.}
\label{tab:lengthclass}
\end{table}

The rank ordering is identical for both length classes (BM25+PRF $>$ BM25 $>$ LSA $>$ ESA-Cranfield $>$ Baseline $>$ QE $>$ ESA-WordNet). PRF's lead over plain BM25 is somewhat more pronounced for medium-length queries (2.9\% gap) than for long queries (2.5\%), which is consistent with the idea that shorter queries give BM25's first pass fewer terms to work with, so the PRF expansion term injection contributes more.

\begin{table}[H]
\centering
\renewcommand{\arraystretch}{1.2}
\begin{tabular}{lccc}
\toprule
\textbf{System} & \textbf{Sparse (1--3 rel.)} & \textbf{Medium (4--9 rel.)} & \textbf{Dense ($\ge$10 rel.)} \\
\midrule
VSM (Baseline)      & 0.4195 & 0.2813 & 0.1680 \\
BM25                & \textbf{0.4820} & 0.3383 & 0.2369 \\
BM25 + PRF          & 0.4251 & \textbf{0.3557} & \textbf{0.2626} \\
LSA ($k{=}200$, CV) & 0.4672 & 0.3275 & 0.2405 \\
ESA (Cranfield)     & 0.4273 & 0.3044 & 0.2132 \\
VSM + QE            & 0.3589 & 0.2488 & 0.1280 \\
ESA (WordNet)       & 0.3224 & 0.2199 & 0.1380 \\
\bottomrule
\end{tabular}
\caption{Mean AP@10 by qrel density class. Bold entries are best per column. All systems score higher on sparse queries because reaching 1--3 relevant documents in top-10 is structurally easier than achieving recall over a 10+ relevant set.}
\label{tab:densityclass}
\end{table}

The density breakdown exposes a leadership inversion: \textbf{BM25} leads the sparse-density class (0.4820) but \textbf{BM25+PRF} leads the medium and dense classes (0.3557 and 0.2626). This matches what we saw in the Marginal Difference Analysis: PRF's expansion step depends on reliable first-pass retrieval, and that becomes more likely when many relevant documents are spread across the collection. The MAP numbers also drop consistently as density increases --- this is a structural property of the metric itself. A query with only 2 relevant documents needs just 2 hits in the top-10 for MAP = 1.0, while a query with 20 relevant documents can at best score around 0.49 at $k{=}10$ under perfect retrieval. So MAP comparisons across density classes need to account for this built-in ceiling difference.

\subsection{Universal Failure Analysis}

A \textbf{universal failure} is a query for which every evaluated system scores AP@10 = 0 --- cases no implemented method can handle. We found 8 such queries out of 225 (3.6\% of the collection).

\begin{table}[H]
\centering
\renewcommand{\arraystretch}{1.2}
\begin{tabular}{cllcc}
\toprule
\textbf{QID} & \textbf{Query text (truncated)} & \textbf{Length} & \textbf{Density} \\
\midrule
9   & papers on internal /slip flow/ heat transfer studies        & medium & medium \\
22  & did anyone else discover that the turbulent skin friction... & long   & sparse \\
28  & what application has the linear theory design of curved...  & long   & sparse \\
44  & what are the details of the rigorous kinetic theory...      & long   & medium \\
62  & how far around a cylinder and under what conditions...      & long   & medium \\
109 & panels subjected to aerodynamic heating                    & medium & medium \\
207 & how do large changes in new mass ratio affect wing-flutter  & long   & medium \\
216 & what investigations have been made of the wave system...    & long   & sparse \\
\bottomrule
\end{tabular}
\caption{Universal failure queries: AP@10 = 0 for all seven systems. Classes from the query-class analysis.}
\label{tab:univfail}
\end{table}

Six of eight are long queries, and three are sparse-density (few relevant documents to find). The noteworthy outlier is Query 109 --- \textit{``panels subjected to aerodynamic heating''} --- which is short (medium length class) and medium density. This query fails despite a reasonable structure because ``panels'' in the Cranfield corpus primarily refers to \emph{structural panels} (thin plates under load) rather than the more common aeronautics sense; without any term that matches the relevant document vocabulary, no method in our set can recover. This shows that vocabulary mismatch --- the same problem that failed the baseline --- is still the root cause even after six different improvements. None of the methods in our set can recover when the query and the relevant documents simply do not share surface vocabulary and the corpus does not have enough co-occurrence signal to bridge the gap.

\subsection{Systematic Failure Analysis}

\begin{table}[H]
\centering
\renewcommand{\arraystretch}{1.2}
\begin{tabular}{lcc}
\toprule
\textbf{System} & \textbf{New Failures (vs.\ Baseline)} & \textbf{Baseline Failures Fixed} \\
\midrule
BM25           & 1  & 7  \\
BM25 + PRF     & 4  & 5  \\
LSA ($k{=}200$, CV) & 5  & 6  \\
ESA (Cranfield)& 8  & 6  \\
ESA (WordNet)  & 16 & 2  \\
VSM + QE       & 16 & 0  \\
\bottomrule
\end{tabular}
\caption{New failure cases introduced vs.\ baseline failures fixed, where a failure is defined as AP@10 = 0. The improved diagnostic in the code also records query length and relevant-doc count for each new failure.}
\end{table}

A ``new failure'' means a query for which the baseline found at least one relevant document in the top-10 but the new method found none. VSM+QE introduces 16 new failures while fixing zero, confirming that query drift is a systematic problem. BM25 is the most conservative: only 1 new failure against 7 fixed --- it changes the most failures in the correct direction while almost never breaking a query that was working.

\subsubsection*{Boundary Conditions: When Do Methods Introduce New Failures?}

The table below breaks down new failures by query length and qrel density to identify the specific conditions under which each method degrades.

\begin{table}[H]
\centering
\renewcommand{\arraystretch}{1.2}
\begin{tabular}{lcc|ccc}
\toprule
\textbf{System} & \multicolumn{2}{c|}{\textbf{By query length}} & \multicolumn{3}{c}{\textbf{By qrel density}} \\
& \textbf{Medium} & \textbf{Long} & \textbf{Sparse} & \textbf{Medium} & \textbf{Dense} \\
\midrule
BM25               & 0 & 1 & 0 & 1 & 0 \\
BM25 + PRF         & 1 & 3 & 1 & 3 & 0 \\
LSA ($k{=}200$)    & 0 & 5 & 0 & 3 & 2 \\
ESA (Cranfield)    & 0 & 8 & 3 & 4 & 1 \\
ESA (WordNet)      & 1 & 15 & 3 & 7 & 6 \\
VSM + QE           & 2 & 14 & 2 & 6 & 8 \\
\bottomrule
\end{tabular}
\caption{New failures broken down by query length and qrel density class. No method introduces failures on short queries because no Cranfield query has $\le 5$ words after preprocessing.}
\label{tab:boundary}
\end{table}

Long queries ($\ge 11$ words) account for the vast majority of new failures across all methods. BM25's single new failure is a 13-word acoustic-propagation query with medium density --- a topic whose vocabulary is far enough from mainstream aerodynamics that BM25's term weighting still falls short. VSM+QE introduces 8 new failures on dense-qrel queries and 6 on medium, confirming that vocabulary drift hurts most when many relevant documents exist and small ranking errors cascade into large recall losses.

\subsection{Qualitative Case Study}

We focus on the documented vocabulary mismatch failure: Query 1 and Document 29.

\textbf{Query:} \textit{``what similarity laws must be obeyed when constructing aeroelastic models of heated high speed aircraft''}

\textbf{Document 29} (relevant, qrel position 2) discusses \textit{``transient temperatures and thermal stresses in aircraft structures subjected to aerodynamic heating''}.

\begin{table}[H]
\centering
\renewcommand{\arraystretch}{1.2}
\begin{tabular}{lc}
\toprule
\textbf{System} & \textbf{Rank of Doc 29} \\
\midrule
VSM (Baseline) & 109 \\
BM25           & 79  \\
LSA ($k{=}200$, CV) & 77  \\
VSM + QE       & 117 (worse than baseline) \\
BM25 + PRF     & 39  \\
ESA (WordNet)  & 126 (worst) \\
ESA (Cranfield)& \textbf{27} (best) \\
\bottomrule
\end{tabular}
\caption{Rank assigned to Document 29 for Query 1 by each system.}
\end{table}

Tracing the rank of Document 29 across all seven systems shows exactly why each method either helps or fails here:

BM25 improves over VSM (rank 79 vs.\ 109) because its Robertson IDF gives moderate-frequency terms like \textit{aircraft} more discriminative weight, and its length normalisation prevents short abstracts from being systematically outscored by longer ones. However, it still cannot bridge the \textit{heated}/\textit{thermal} lexical gap.

LSA (rank 77) provides a partial but incomplete correction. The $k{=}200$ latent space captures some of the co-occurrence signal between heat-related terms, but the aerodynamics corpus is narrow enough that even 200 components do not fully separate the thermal cluster.

BM25+PRF (rank 39) achieves the second-best result. The first-pass BM25 retrieves documents that use \textit{thermal} and \textit{aerodynamic}, and the PRF expansion injects these terms into the second-pass query, effectively bridging the lexical gap through an unsupervised expansion step.

ESA-Cranfield (rank 27) is the best performer for this specific query. By projecting both the query and documents into the full Cranfield document-level concept space, it captures the fact that Document 29 has a high similarity profile with documents that are themselves relevant to heated aerodynamic structures.

VSM+QE (rank 117, worse than baseline) exemplifies query drift. The primary WordNet synset for \textit{heated} or \textit{aeroelastic} likely retrieves semantically shallow synonyms that actively mislead the TF-IDF ranking, pushing Document 29 even further down.

\subsection{Discussion: Why the Methods Work (or Fail)}

The main pattern across all results is whether a method's assumptions actually fit the Cranfield domain. Here is a per-method breakdown with supporting data:

\begin{itemize}

\item \textbf{BM25 worked} because TF saturation and length normalisation fix two concrete problems the VSM has on this corpus. TF saturation stops high-frequency terms from dominating (a term appearing 50 times is not treated as $50\times$ more relevant than one appearing once). Length normalisation stops long abstracts from outscoring short focused ones unfairly. The CV selected $k_1{=}2.0$ and $b{=}0.9$ --- both higher than standard defaults of 1.5/0.75 --- showing that Cranfield needs stronger normalisation due to its high variance in abstract length. \textit{Proof:} MAP jumped from 0.2691 to 0.3304 (+22.8\%). BM25 improved 142 out of 225 queries and degraded only 37 --- a 3.8:1 ratio in favour of improvement.

\item \textbf{LSA worked} because the Cranfield corpus is a single narrow domain (aerodynamics and fluid mechanics), so the SVD at $k{=}200$ can extract consistent co-occurrence structure. Terms like \textit{heated} and \textit{thermal} end up with nearby latent representations because they co-occur with the same pool of aerodynamics terms across the 1400 documents. \textit{Proof:} MAP at $k{=}200$ (0.3233) beats $k{=}300$ (0.3205) in the ablation --- going beyond 200 starts pulling in noisy dimensions and actually hurts. CV independently confirmed $k{=}200$ without seeing test labels.

\item \textbf{BM25 + PRF worked for dense queries but not for sparse ones.} PRF assumes the first-pass top-3 documents are relevant enough to yield useful expansion terms. When many relevant documents exist, BM25 reliably surfaces at least one in the top-3 and the expansion terms genuinely reflect the query topic. When the relevant set is small, the top-3 may all be off-topic and the second pass gets misled. \textit{Proof:} PRF wins 43 of 67 dense-qrel queries (64.2\%) but only 5 of 35 sparse-qrel queries (14.3\%). The 0.83\% aggregate MAP advantage for PRF exists because Cranfield has 67 dense queries vs.\ only 35 sparse ones --- the distribution tips the average in PRF's favour.

\item \textbf{WordNet Query Expansion hurt performance} because Cranfield uses highly specialised 1950s aerodynamics vocabulary. Terms like \textit{aeroelastic}, \textit{blasius equation}, or \textit{joule heating} either do not appear in WordNet at all, or map to general-English synsets that are useless in this domain. Expanding \textit{viscous} with everyday synonyms for \textit{thick} or \textit{sticky} adds noise that pushes relevant documents down. \textit{Proof:} QE degraded 130 queries while improving only 28 (a 4.6:1 ratio \emph{against} improvement). The two-sided Wilcoxon test gives $p = 8.1 \times 10^{-13}$, $\bar{\delta} = -0.039$ MAP --- statistically significantly worse than baseline.

\item \textbf{ESA with WordNet concepts failed} for the same reason as QE: the concept space was built for everyday English, not for a 1950s aerodynamics corpus. Terms central to Cranfield like \textit{laminar}, \textit{stagnation point}, or \textit{oblique shock} are either missing from WordNet or too weakly connected to form useful concept dimensions. \textit{Proof:} ESA-WordNet degraded 120 queries while improving only 60. MAP fell to 0.2115 (below baseline 0.2691), $p = 4.0 \times 10^{-9}$ --- also significantly worse.

\item \textbf{ESA with Cranfield concepts worked better than WordNet ESA} because the concept space is built from the same 1400 documents, so every specialised aerodynamics term gets a concept vector that reflects its actual usage in context. Domain mismatch is avoided by construction. It is still not the best method overall because projecting into a 1400-dimensional concept space is noisy and high-dimensional. \textit{Proof:} MAP 0.2963 vs.\ 0.2115 for WordNet ESA. It improved 120 queries and degraded 71, a meaningful improvement over WordNet ESA's 60/120 profile.

\end{itemize}

\subsection{Summary of Findings}

\begin{enumerate}[noitemsep]
    \item All six methods are statistically significantly different from the VSM baseline (Wilcoxon, $p < 0.05$). Four improve it; two degrade it.
    \item BM25 ($k_1{=}2.0$, $b{=}0.9$, CV-selected) and BM25+PRF are the most practically reliable: highest per-query improvement rates (142 and 136 improved queries), fewest new failures (1 and 4), and strong aggregate metrics.
    \item The MAP gap between BM25 (0.3304), BM25+PRF (0.3388), and LSA (0.3233) is marginal ($<1.6\%$) and does not constitute a credible superiority claim without pair-level tests. Per-query analysis shows BM25+PRF wins on more individual queries than BM25.
    \item VSM+QE and ESA-WordNet are statistically significantly \emph{worse} than the baseline, not merely indistinguishable --- a distinction only visible when using the correct two-sided Wilcoxon test.
    \item The number of LSA latent dimensions is best selected by 5-fold CV ($k{=}200$, CV MAP = 0.3233). Selecting $k$ on all 225 test queries constitutes data leakage; CV gave the same answer here, but the methodology matters.
    \item MRR and graded nDCG provide complementary insight: BM25 leads on first-relevant-position quality (MRR = 0.8002); BM25+PRF leads on surfacing the most highly relevant documents (graded nDCG = 0.4887).
    \item The vocabulary mismatch case study confirms that all three corpus-based methods (BM25, LSA, PRF) partially address the core failure, but none fully resolves it.
\end{enumerate}

% ===========================================================================
\section{Conclusion}
% ===========================================================================

This project built and evaluated an Information Retrieval system on the Cranfield aerodynamics dataset, starting from a TF-IDF Vector Space Model baseline and testing six alternative retrieval approaches. The baseline VSM achieved MAP@10 = 0.2691 and nDCG@10 = 0.4203. The main problem identified was the orthogonality assumption, which caused vocabulary mismatch failures like Document 29 being placed at rank 109 for Query 1 despite being clearly relevant.

Six alternative systems were tested. BM25 ($k_1{=}2.0$, $b{=}0.9$, hyperparameters selected by 5-fold cross-validation) gave the most consistent improvement: MAP@10 = 0.3304 (+22.8\% relative), Wilcoxon $p = 3.3 \times 10^{-16}$, Cohen's $d = 0.57$ (medium effect), improving 142 out of 225 queries and degrading only 37. BM25+PRF achieved MAP@10 = 0.3388 with the best graded nDCG (0.4887), and LSA at $k{=}200$ (selected by 5-fold CV to avoid data leakage) achieved MAP@10 = 0.3233. All three are statistically significantly better than the baseline.

Importantly, the correct application of two-sided Wilcoxon tests for VSM+QE and ESA-WordNet revealed that they are not merely indistinguishable from the baseline --- they are statistically significantly \emph{worse} ($p < 10^{-8}$). This distinction matters: one-sided tests had previously masked the harm these methods cause.

The overall takeaway is that methods which rely on the corpus itself (BM25's statistics, LSA's SVD, PRF's feedback) worked, while methods that import outside knowledge from WordNet failed because the domain mismatch was too large. One direction to explore further would be transformer-based dense retrieval, which learns semantic representations directly from text rather than depending on a hand-built lexical resource like WordNet.

\end{document}
